\section{Related Work and Positioning}

\textbf{Inference-energy measurement.}  Inference energy depends on token
geometry, model, framework, and hardware
\cite{samsi2023words,fernandez2025energy}.  ML.ENERGY automates measurement
across systems, while \bench{} contributes
declarative workloads and phase-aligned device/node telemetry
\cite{chung2025mlenergy,niu2026tokenpowerbench}.  These systems answer what a
completed run consumed.  \system{} retains real-GPU measurement as calibration
source and final authority, while CPU-resident simulated or extrapolated
records estimate which runs to purchase next.  Its phase boundary preserves the
semantics needed to compare them.

\textbf{Serving simulation and recommendation.}  Vidur provides event-driven
performance simulation for LLM serving; LLM-Pilot characterizes models and
hardware to recommend deployments; and Charon models fine-grained large-scale
training and inference \cite{agrawal2024vidur,lazuka2024llmpilot,
yang2026charon}.  Our current model is deliberately local: an analytical
decomposition supplies a prior, six
measured workloads learn only a sparse residual, and an exact configuration
contract bounds reuse.  It neither creates virtual GPUs nor claims that one
H100 measurement reproduces communication, queueing, or power behavior on a
multi-node target.  Its differentiator is the release boundary: split labels,
frozen hashes, abstention, and an independent run are first-class outputs, not
evaluation bookkeeping added after a recommendation.

\textbf{Energy-aware serving optimization.}  DynamoLLM designs inference
clusters under performance and energy objectives; BEAM jointly tunes resources
and power under SLOs; and BOute applies multi-objective Bayesian optimization to
heterogeneous LLM serving \cite{stojkovic2025dynamollm,lee2026beam,
jiang2026boute}.  These systems are closest to the eventual deployment use
case, but solve a different layer: a policy searches a configuration space
given an evaluator.  This WIP audits the evaluator itself.  It therefore
reports blind point error, ordering, false-confidence risks, and purchased
measurement cost before claiming energy savings.  A future TokenPowerAgent can
compare random, Bayesian, and agentic policies over the same evidence API; no
policy will be allowed to promote a predicted record to measured evidence or
train on its final holdout.  The present result therefore supports same-GPU
workload screening within a measured envelope; configuration, topology, and
hardware transfer remain separate experiments in the claim ladder.
